\documentclass[10pt,aspectratio=169,dvipsnames]{beamer}
\usetheme[]{Berlin}

\setbeamertemplate{footline}{
  \usebeamercolor[fg]{framesource}%
  \usebeamerfont{page number in head}%
  \hspace{0.2cm}
  \small \insertframenumber
  \vspace{0.2cm}
  \doclicenseIcon
  \hfill
  
\includegraphics[height=0.8cm]{logos/tub_logo.pdf}
  \hspace{0.2cm}
}

\setbeamercovered{transparent}

\setbeamertemplate{footline}[
 myframe number]

% PACKAGES
\usepackage[absolute,overlay]{textpos}
\usepackage[utf8]{inputenc}
\usepackage[official]{eurosym}
\usepackage{booktabs,multicol,multirow,tabularx,array}
\usepackage{subcaption}
\usepackage{parskip}
\usepackage{bm}
\usepackage{tikz}
\usepackage{adjustbox}
\usepackage[super]{nth}
\usepackage[version=4]{mhchem}
\usepackage{siunitx}

\sisetup{
  detect-all = true
}


\usepackage[
    type={CC},
    modifier={by},
    version={4.0},
]{doclicense}

% HYPERREFERENCES
\usepackage{hyperref}
\hypersetup{
	colorlinks=true,
	citecolor=tub-blue,
	linkcolor=tub-blue,
	urlcolor=tub-blue
}

\usepackage[sfdefault]{roboto}
\usepackage[normalem]{ulem}

% \usepackage[backend=biber,style=numeric-comp]{biblatex}
% \addbibresource{references.bib}

\newcolumntype{R}[1]{>{\raggedleft\arraybackslash}p{#1}}


% GRAPHICS
\graphicspath{
    {graphics/},
    {graphics/logos/},
  }
\DeclareGraphicsExtensions{.pdf,.jpeg,.png,.jpg}

% FORMATTING
\setlength{\parskip}{6pt}
\linespread{1.1}
\newcommand{\seprule}{\par\noindent\textcolor{black!25}{\rule{\textwidth}{0.4pt}}}

% SOURCES
\setbeamercolor{framesource}{fg=gray}
\setbeamerfont{framesource}{size=\tiny}
\newcommand{\source}[1]{\begin{textblock*}{10cm}(3.6cm,8.25cm)
    \begin{beamercolorbox}[ht=0.5cm,right]{framesource}
        \usebeamerfont{framesource}\usebeamercolor[fg]{framesource} Quelle: {#1}
    \end{beamercolorbox}
\end{textblock*}}

% REFERENCES
% \usepackage[style=nature, doi=true, maxcitenames=3]{biblatex}
% \bibliography{bibliography}
% \renewcommand*{\bibfont}{\footnotesize}

% COLOR TEXT
\newcommand{\bl}[1]{\textcolor{tub-blue}{#1}}
\newcommand{\gr}[1]{\textcolor{tub-green}{#1}}
\newcommand{\rd}[1]{\textcolor{tub-red}{#1}}
\newcommand{\yl}[1]{\textcolor{tub-yellow}{#1}}
% \newcommand{\org}[1]{\textcolor{tub-orange}{#1}}
\newcommand{\cb}[1]{\colorbox{gray!20}{#1}}

\usepackage{array}% https://ctan.org/pkg/array
\makeatletter
\g@addto@macro{\endtabular}{\rowfont{}}% Clear row font
\makeatother
\newcommand{\rowfonttype}{}% Current row font
\newcommand{\rowfont}[1]{% Set current row font
   \gdef\rowfonttype{#1}#1%
}
\newcolumntype{L}{>{\rowfonttype}l}

\usepackage[most]{tcolorbox}
\tcbuselibrary{listingsutf8}

% Custom commands
\usepackage{caption}
\renewcommand{\figurename}{} % Removes the "Figure" text
\captionsetup[figure]{labelformat=empty, labelsep=none} % Removes numbering and colon


% FONT
%\usepackage{utopia}

% REFERENCES
%\usepackage[backend=biber,style=authoryear-comp]{biblatex}
%\bibliography{references.bib}

% TITLE PAGE
\title{
  \textbf{IN4climate --- Kickoff Sprint CO$_2$ Infrastruktur}
}
\vspace*{-.5cm}

% \author{STRIse}
\institute[Technische Universität Berlin] % (optional, but mostly needed)
{
  \normalsize
  \alert{Kurzvorstellung des RESILIENT-Projekts} \\
  \footnotesize
  \url{https://resilient-project.github.io/} \\
  FG Digitaler Wander in Energiesystemen \\
  Technische Universität Berlin \\

  \vspace{0.3cm}

  07. Oktober 2025
}

\date{}

\subject{IN4climate --- Kickoff Sprint CO2 Infrastruktur}

\titlegraphic{%
\vspace{-1cm}

\includegraphics[height=1cm,clip=true]{logos/cetp_logo.png}
\hspace{0.5cm}
\includegraphics[height=1cm,clip=true]{logos/ensys_short_logo.pdf}
\hspace{0.5cm}
\includegraphics[trim=0 0cm 0 0cm,height=1cm,clip=true]{logos/tub_logo.pdf}
\hspace{0.5cm}\includegraphics[trim=0.2cm 0.6cm 0.6cm 0.2cm,height=1.1cm,clip=true]{logos/bmwk_en_logo.png}
}

\begin{document}

\addtocounter{framenumber}{-1}
{
	\setbeamertemplate{footline}{
		\usebeamercolor[fg]{framesource}%
		\usebeamerfont{page number in head}%
		\begin{center}
			\Large
			\doclicenseIcon
			\vspace{0.4cm}
		\end{center}
	}
	\maketitle
}

\section{Projektvorstellung}
\begin{frame}{RESILIENT Projektverbund}
	\centering
	\includegraphics[width=0.85\textwidth]{other/resilient_partners}

	\footnotesize
	Förderung über das \alert{CETPartnership 2022} Programm, sowie \alert{BMWK} für die deutschen Projektpartner.

	\source{\url{https://resilient-project.github.io/}}
\end{frame}

\begin{frame}{RESILIENT-Projektübersicht}
	\footnotesize
	\begin{columns}
		\begin{column}{0.75\textwidth}
			\begin{itemize}
				\item \alert{Problemdiagnose}: Welt befindet sich in einer Phase der \alert{Unsicherheit} und des \alert{Wandels}
				      (Geopolitik, Energiewende, Klimawandel); Szenarien mit einzelnen deterministischen
				      Lösungen nicht mehr bedarfsgerecht
				\item \alert{Projektziel}: Entwicklung und Demonstration von Werkzeugen für eine resiliente Planung der Energiewende für Industrie, E-fuels und Netze
				\item Baut auf dem schon etablierten, quelloffenen und weltweit genutzten Multi-Vektor-Energiesystemmodell \alert{Python for Power System Analysis (PyPSA)}
				\item \alert{Methodischer Schwerpunkt}: Optimierung unter Unsicherheit (stochastisch und robust),  Dekompositionsmethoden und Heuristiken für Resilienz
				\item \alert{Fallstudien}: \alert{Industriewende in NRW}, Netzplanung in BW, E-Fuels in Skandinavien, Klimaextremereignisse in Frankreich
			\end{itemize}
		\end{column}
		\begin{column}{0.25\textwidth}
			
\includegraphics[width=1\textwidth]{logos/logo-secondary-light.png}
		\end{column}
	\end{columns}

\end{frame}

\begin{frame}{Arbeitspakete}
	\centering
	\includegraphics[width=0.8\textwidth]{other/resilient_project_structure}
\end{frame}

\begin{frame}{PyPSA-Eur: Ein sektorenübergreifendes Energiesystemmodell}
	\begin{columns}
		\begin{column}{0.5\textwidth}
			\footnotesize
			\begin{itemize}
				\setlength\itemsep{.8em}
				\item Räumlich und zeitlich hochaufgelöstes lineares Optimierungsmodell, das den \alert{europäischen} Kontinent abdeckt,
				\item Aufbauend auf dem Open-Source-Framework \alert{PyPSA},
				\item Berücksichtigt den \alert{Bestand} bestehender Kraftwerke, erneuerbare Potenziale, \alert{Verfügbarkeitszeitreihen},
				\item Umfasst das \alert{elektrische Übertragungsnetz} von AC 220 kV bis 750 kV (UA) sowie DC ab 150 kV, mit der Option, geplante Netzausbauprojekte (TYNDP und deutscher NEP) einzubeziehen, ebenso das Gasnetz,
				\item Betreuung und Entwicklung durch das FG Digitaler Wandel in Energiesystemen an der \alert{TU Berlin}.
			\end{itemize}
		\end{column}
		\begin{column}{0.5\textwidth}
			\includegraphics[width=1\textwidth]{other/pypsa_eur_elec}
		\end{column}
	\end{columns}

	\source{\url{https://pypsa-eur.readthedocs.io}}
\end{frame}

% \begin{frame}{Selection of planned model developments}
%   \footnotesize
%   \begin{columns}[T] % align columns
%     \begin{column}{.5\textwidth}
%         \begin{minipage}[t][.45\textheight]{\linewidth}
%             \begin{alertblock}{\textbf{Computational methods for uncertainties}}
%                 \begin{itemize}
%                   \item decomposition techniques
%                   \item large-scale stochastic optimisation
%                   \item \textbf{test robustness of system}
%                   \item using SMS++ framework
%                 \end{itemize}
%             \end{alertblock}
%         \end{minipage}

%         \vspace{-0.03\textheight} % separation between the boxes

%         \begin{minipage}[t][.45\textheight]{\linewidth}
%             \begin{alertblock}{\textbf{Industry transformation (FORECAST)}}
%               \begin{itemize}
%                 \item fuel and process switching
%                 \item industry relocation
%                 \item carbon sources and feedstocks
%                 \item data on stock \& investment cycles
%                 \item new technologies (oxyfuel cement, etc.)
%               \end{itemize}
%             \end{alertblock}
%         \end{minipage}
%     \end{column}

%     \begin{column}{.5\textwidth}
%         \begin{minipage}[t][.45\textheight]{\linewidth}
%             \begin{alertblock}{\textbf{Carbon management and biomass usage}}
%               \begin{itemize}
%                 \item \textbf{CO$_2$ network}
%                 \item \textbf{CO$_2$ sequestration potentials}
%                 \item circular carbon economy and recycling
%                 \item biomass usage options
%               \end{itemize}
%             \end{alertblock}
%         \end{minipage}

%         \vspace{-0.03\textheight} % separation between the boxes

%         \begin{minipage}[t][.45\textheight]{\linewidth}
%             \begin{alertblock}{\textbf{Global green fuel and material markets}}
%               \begin{itemize}
%                 \item \textbf{imports of green energy and materials}
%                 \item \textbf{effects on European infrastructure}
%                 \item restructuring of value chains
%                 \item risks (geopolitical, technological, etc.) \newline
%               \end{itemize}
%             \end{alertblock}
%         \end{minipage}
%     \end{column}
% \end{columns}
% \end{frame}

\section{Aktuelle Studienbeispiele}

\begin{frame}{Studie: H$_2$ zu CO$_2$ oder CO$_2$ zu H$_2$ Standorten transportieren?}

	\centering
	\includegraphics[width=0.35\textwidth]{other/carbon_networks_balance_map_carbon}
	\includegraphics[width=0.35\textwidth]{other/carbon_networks_balance_map_hydrogen}

	\source{Hofmann, Tries, Neumann, Zeyen, Brown (2025): \\\url{https://www.nature.com/articles/s41560-025-01752-6}}
\end{frame}

\begin{frame}{Studie: PCI-PMI H$_2$- und CO$_2$-Infrastrukturprojekte}
	\scriptsize
	\begin{columns}[T] % align columns
		% Erste Spalte
		\begin{column}{.63\textwidth}
			\begin{minipage}[t][.45\textheight]{\linewidth}
				\begin{alertblock}{\textbf{Was sind PCI-PMI-Projekte?}}
					\begin{itemize}
						\setlength\itemsep{0.85em}
						\item Projects of Common Interest (PCI) sind zentrale \alert{grenzüberschreitende Infrastrukturprojekte}, die die Energiesysteme der EU-Mitgliedsstaaten verbinden
						\item Ihr Ziel ist es, `der EU beim Erreichen ihrer \alert{energiepolitischen und klimapolitischen Ziele} zu unterstützen: bezahlbare, sichere und nachhaltige Energie für alle Bürgerinnen und Bürger sowie die langfristige Dekarbonisierung der Wirtschaft im Einklang mit dem \alert{Pariser Abkommen}'
						\item Der potenzielle Gesamtnutzen jedes Projekts muss die Kosten übersteigen.
						\item Aufgrund ihres \alert{Leuchtturmcharakters} ist die Umsetzung dieser Projekte sehr wahrscheinlich
						\item Große Infrastrukturprojekte (einschließlich PCI-PMI) sind jedoch häufig von Verzögerungen durch Genehmigungen, Engpässe bei der Beschaffung usw. betroffen
					\end{itemize}
				\end{alertblock}
			\end{minipage}
		\end{column}

		% Second column
		\begin{column}{.36\textwidth}
			\begin{minipage}[t][.45\textheight]{\linewidth}
				\begin{alertblock}{\textbf{Project map}}
					\centering
					\includegraphics[width=1\textwidth]{map_adm_pcipmi}
				\end{alertblock}
			\end{minipage}
		\end{column}

	\end{columns}

	\vspace{2.1cm}

	\begin{enumerate}
		\item Welchen Beitrag erbringen PCI-PMI-Projekte bei der Erreichung der EU \\Klima- und Energiepolitikziele und welche Kosten sind damit verbunden?
		\item Welche Kosten entstehen bei der Einhaltung der Ziele bei einer Verzögerunng dieser Projekte?
	\end{enumerate}

	\source{Eigene Darstellung nach \url{https://ec.europa.eu/energy/infrastructure/transparency_platform/map-viewer}. Preprint Xiong, Brown, Riepin (2025): \url{https://doi.org/10.48550/arXiv.2507.01860}}

\end{frame}

\begin{frame}{Studie: PCI-PMI H$_2$- und CO$_2$-Infrastrukturprojekte}
	\centering
	\begin{columns}
		\begin{column}{0.5\textwidth}
			\footnotesize
			\begin{figure}
				\centering
				\includegraphics[width=0.8\textwidth]{pci___base_s_adm___2050-balance_map_H2}
			\end{figure}
		\end{column}
		\begin{column}{0.5\textwidth}
			\begin{figure}
				\centering
				\includegraphics[width=0.8\textwidth]{pci___base_s_adm___2050-balance_map_co2_stored}
			\end{figure}
		\end{column}
	\end{columns}
	\source{Preprint Xiong, Brown, Riepin (2025): \url{https://doi.org/10.48550/arXiv.2507.01860}}
\end{frame}

\begin{frame}
	\setbeamercolor{background canvas}{bg=red}
	\centering
	\fcolorbox{white}{tub-red}{\parbox[c][0.9\textheight][c]{\textwidth}{%
			\centering \color{white}\Large{\textbf{Vielen Dank.}} \\
			\scriptsize
			\vspace{0.5cm}
			\textbf{RESILIENT} \\
			\textit{Resilient Energy System Infrastructure Layouts for Industry, E-Fuels and Network Transitions} \\
			\href{https://resilient-project.github.io}{\textcolor{white}{$\hookrightarrow$ resilient-project.github.io}} \\
			\href{https://github.com/pypsa/pypsa-eur}{\textcolor{white}{$\hookrightarrow$ github.com/pypsa/pypsa-eur}} \\
			\vspace{0.5cm}
			Department of\\\textbf{Digital Transformation in Energy Systems (ENSYS)} \\
			\vspace{0.5cm}
			\textbf{Projektleitung: Dr. Iegor Riepin}\\
			\href{mailto:iegor.riepin@tu-berlin.de}{\textcolor{white}{$\hookrightarrow$ iegor.riepin@tu-berlin.de}} \\
			\vspace{0.5cm}
			\textbf{Präsentation: Bobby Xiong}\\
			\href{mailto:xiong@tu-berlin.de}{\textcolor{white}{$\hookrightarrow$ xiong@tu-berlin.de}} \\
		}}
\end{frame}

% Appendix
\section{Appendix}

\begin{frame}{PyPSA-Eur: Eine typische Modellkonfiguration}
	\begin{columns}
		\begin{column}{0.5\textwidth}
			\footnotesize
			\begin{itemize}
				\setlength\itemsep{.8em}
				\item \alert{Alle Energiesektoren} (Strom, Wärme, Verkehr, Industrie, Rohstroffe, Landwirtschaft)
				\item \alert{Infrastrukturplanung} von Erzeugung, Übertragung, Speicherung und Power-to-X mit \alert{Technologieprognosen}
				\item \alert{100-200 Regionen} in European
				\item \alert{2-3 stündliche} segmentierte Zeitreihen
				\item \alert{Netto-Null-Emissionen} bis 2050
				\item Minimiere System in \alert{Optimierung}
				\item Viele \alert{Szenarien} durchspielen (oder \alert{stochastische Optimierung})
			\end{itemize}
		\end{column}
		\begin{column}{0.5\textwidth}
			\includegraphics[width=1\textwidth]{other/multisector_diagram}
		\end{column}
	\end{columns}

	\source{\url{https://pypsa-eur.readthedocs.io}}
\end{frame}

\begin{frame}{Annahmen zur Nachfrageverteilung pro Sektor}
	\centering
	\begin{columns}
		\begin{column}{0.48\textwidth}
			\footnotesize
			\begin{figure}
				\centering
				\includegraphics[width=1\textwidth]{exogenous_demand}
			\end{figure}
		\end{column}
		\begin{column}{0.52\textwidth}
			\footnotesize
			\begin{itemize}
				\item Nachfrage nach Strom, Wärme, Gas, Biomasse und Verkehr ist räumlich und zeitlich aufgelöst
				\item Verbrennungsmotoren im Landverkehr werden voraussichtlich bis 2050 vollständig zugunsten von Elektrofahrzeugen abgelöst
				\item Nachfrage nach Methanol und Kohlenwasserstoffen, einschließlich Kerosin, wird hauptsächlich durch Schifffahrt, Luftfahrt und Industrie getrieben (nicht räumlich aufgelöst)
				\item Nicht vermeidbare Prozessemissionen aus der Industrie, z.\,B. Zement, werden ebenfalls berücksichtigt
			\end{itemize}
		\end{column}
	\end{columns}
	\source{Preprint Xiong, Brown, Riepin (2025): \url{https://doi.org/10.48550/arXiv.2507.01860}}
\end{frame}

\begin{frame}{Studie: Nutzenvergleich Stromnetzausbau vs. H$_2$-Netz}
	\centering
	\includegraphics[width=0.9\textwidth]{other/neumann_hydrogen_paper}
	\\
	\footnotesize
	$\rightarrow$ Wasserstoffnetz \alert{reduziert Systemkosten}, aber \alert{Stromnetzausbau ist attraktiver.}

	\source{Neumann, Zeyen, Victoria, Brown: \url{https://doi.org/10.1016/j.joule.2023.06.016}}
\end{frame}

\begin{frame}{H$_2$-Nutzung hauptsächlich für Derivate und als Grundstoff}

	\includegraphics[width=1.2\textwidth, trim=0cm 6.6cm 0cm 6.6cm, clip]{other/hydrogen_balance.pdf}

	\footnotesize

	\alert{Wenig direkte Anwendungen von Wasserstoff}; meist zur Synthese anderer Derivate

	\begin{multicols}{2}
		\begin{itemize}
			\item Ammoniak für Düngemittel
			\item Direktreduktion von Eisen für Stahl
			\item Schiffs- und Flugkraftstoffe
			\item Vorstufen für hochwertige Chemikalien
			\item Backup für Wärme- und Stromversorgung
			\item Einsatz evtl. im Schwerlastverkehr
		\end{itemize}
	\end{multicols}

	\source{Neumann, Zeyen, Victoria, Brown (2023): \\\url{https://doi.org/10.1016/j.joule.2023.06.016}}
\end{frame}

\begin{frame}{CO$_2$: Abscheidung, Nutzung, Transport und Speicherung}

	\begin{center}
		\includegraphics[height=0.73\textheight]{other/carbon_networks_cost_bar}
		\includegraphics[height=0.71\textheight]{other/carbon_networks_balance_bar_carbon}

		\vspace{-0.25cm}
		\footnotesize
		\begin{itemize}
			\item \alert{CCS} für Prozessemissionen (z.\,B. in der Zementindustrie)
			\item \alert{CCU} für E-Synfuels und E-Chemikalien (insbesondere Schifffahrt, Luftfahrt, Kunststoffe)
			\item \alert{CDR} für nicht vermeidbare und negative Emissionen (zum Ausgleich unvollständiger Abscheideraten)
		\end{itemize}

	\end{center}

	\source{Hofmann, Tries, Neumann, Zeyen, Brown (2024): \\\url{https://www.nature.com/articles/s41560-025-01752-6}}

\end{frame}

\begin{frame}{Übertragungsnetz basierend auf OpenStreetMap (OSM)}

	\begin{columns}
		\begin{column}{0.66\textwidth}
			\footnotesize
			\begin{itemize}
				\setlength\itemsep{.8em}
				\item Datensatz enthält das europäische Hoch- und Höchstspannungsnetz (220 kV bis 750 kV), erstellt auf Basis von OSM-Daten
				\item Ein heuristischer Aufbereitungsprozess wurde angewendet für Einträge, bei denen elektrische Parameter unvollständig, fehlend oder mehrdeutig sind
				\item Umspannwerke innerhalb eines Radius von 500 m werden zu gemeinsamen Knoten aggregiert
				\item Für jedes Spannungspaar in einem Umspannwerk werden Transformatoren hinzugefügt
				\item Wechselstromleitungen werden über Standarddaten der pandapower-Bibliothek parametriert und mit einem Faktor von 0.7 multipliziert, um eine n-1-Sicherheitsnäherung zu berücksichtigen
				\item Enthält alle 38 europäischen HGÜ-Verbindungen in ihrer Nennleistung, die bis Ende 2024 in Betrieb genommen wurden
			\end{itemize}
		\end{column}
		\begin{column}{0.34\textwidth}
			\includegraphics[width=1\textwidth]{osm_map}
		\end{column}
	\end{columns}
	\source{Eigene Darstellung auf Basis von OSM-Daten, extrahiert mit der Overpass turbo API\\\url{https://openstreetmap.org}}

\end{frame}

\end{document}
